%=========== ΛΥΣΗ - 2ο ΔΙΑΓΩΝΙΣΜΑ ΠΡΟΣΟΜΟΙΩΣΗΣ ===========
\begin{thema}{Α}
\begin{erwthma}
\item Επειδή $f'(x)>0$ στο $(a,x_0)$, η $f$ είναι γνησίως αύξουσα στο $(a,x_0]$, αφού είναι συνεχής στο $x_0$. Άρα, για κάθε $x\in(a,x_0)$ ισχύει
\[
f(x)<f(x_0).
\]
Επίσης, επειδή $f'(x)<0$ στο $(x_0,\beta)$, η $f$ είναι γνησίως φθίνουσα στο $[x_0,\beta)$, αφού είναι συνεχής στο $x_0$. Άρα, για κάθε $x\in(x_0,\beta)$ ισχύει
\[
f(x)<f(x_0).
\]
Επομένως, για κάθε $x\in(a,\beta)$ με $x\neq x_0$ έχουμε
\[
f(x)<f(x_0).
\]
Άρα το $f(x_0)$ είναι \textbf{τοπικό μέγιστο} της $f$.

\item Μια συνάρτηση $f$ λέγεται γνησίως αύξουσα σε ένα διάστημα $\varDelta$, όταν για οποιαδήποτε $x_1,x_2\in\varDelta$ ισχύει:
\[
x_1<x_2\Rightarrow f(x_1)<f(x_2).
\]

\item Μια συνάρτηση $F$ λέγεται αρχική συνάρτηση ή παράγουσα της $f$ στο $\varDelta$, όταν είναι παραγωγίσιμη στο $\varDelta$ και ισχύει
\[
F'(x)=f(x)
\]
για κάθε $x\in\varDelta$.

\item Οι χαρακτηρισμοί είναι:
\begin{alist}
\item \textbf{Σωστό}. Είναι η ιδιότητα διατήρησης προσήμου συνεχούς συνάρτησης.
\item \textbf{Σωστό}. Από το Θεμελιώδες Θεώρημα του Ολοκληρωτικού Λογισμού,
\[
\dint_a^\beta f'(x)\,\d x=f(\beta)-f(a)=0.
\]
\item \textbf{Σωστό}. Για παράδειγμα, η $f(x)=\dfrac{1}{x}$ στο $\mathbb{R}^*$ είναι $1-1$, αλλά δεν είναι γνησίως μονότονη σε όλο το πεδίο ορισμού της.
\item \textbf{Λάθος}. Μια κατακόρυφη ευθεία τέμνει τη γραφική παράσταση μιας συνάρτησης το πολύ σε ένα σημείο, και τη τέμνει μόνο αν η τετμημένη της ανήκει στο πεδίο ορισμού.
\item \textbf{Σωστό}. Αν $f''(x)>0$ για κάθε $x\in\mathbb{R}$, τότε η $f$ είναι κυρτή σε όλο το $\mathbb{R}$ και δεν αλλάζει κυρτότητα.
\end{alist}
\end{erwthma}
\end{thema}

\begin{thema}{Β}
Δίνονται οι συναρτήσεις
\[
f(x)=\frac{e^x}{e^x+1}
\quad\text{και}\quad
g(x)=\ln x.
\]
\begin{erwthma}
\item Για να ορίζεται η $h=f\circ g$, πρέπει να ορίζεται η $g$, δηλαδή
\[
x>0.
\]
Επειδή $D_f=\mathbb{R}$, δεν υπάρχει άλλος περιορισμός. Άρα
\[
D_h=(0,+\infty).
\]
Για κάθε $x>0$ είναι
\[
h(x)=f(g(x))=f(\ln x)=\frac{e^{\ln x}}{e^{\ln x}+1}
=\frac{x}{x+1}.
\]
Επομένως
\[
{h(x)=\frac{x}{x+1},\quad x>0}.
\]

\item Για κάθε $x>0$ είναι
\[
h'(x)=\left(\frac{x}{x+1}\right)'=\frac{x+1-x}{(x+1)^2}
=\frac{1}{(x+1)^2}>0.
\]
Άρα η $h$ είναι \textbf{γνησίως αύξουσα} στο $(0,+\infty)$, οπότε είναι $1-1$ και αντιστρέφεται.

Για να βρούμε την αντίστροφη, θέτουμε
\[
y=\frac{x}{x+1}.
\]
Τότε
\[
y(x+1)=x
\Leftrightarrow yx+y=x
\Leftrightarrow y=x(1-y)
\Leftrightarrow x=\frac{y}{1-y}.
\]
Επίσης
\[
\lim_{x\to0^+}h(x)=0
\quad\text{και}\quad
\lim_{x\to+\infty}h(x)=1.
\]
Αφού η $h$ είναι γνησίως αύξουσα, έχει σύνολο τιμών
\[
h\bigl((0,+\infty)\bigr)=(0,1).
\]
Άρα
\[
{h^{-1}(x)=\frac{x}{1-x},\quad x\in(0,1)}.
\]

\item Υπολογίζουμε:
\[
\lim_{x\to0^+}h(x)=\lim_{x\to0^+}\frac{x}{x+1}=0
\]
και
\[
\lim_{x\to+\infty}h(x)
=\lim_{x\to+\infty}\frac{x}{x+1}
=\lim_{x\to+\infty}\frac{1}{1+\frac{1}{x}}=1.
\]

\item Θεωρούμε τη συνάρτηση
\[
\varphi(x)=h(x)-e^{-x}=\frac{x}{x+1}-e^{-x},\quad x\in[0,1].
\]
Η $\varphi$ είναι συνεχής στο $[0,1]$ και
\[
\varphi(0)=0-1=-1<0.
\]
Επίσης
\[
\varphi(1)=\frac{1}{2}-\frac{1}{e}>0,
\]
αφού $e>2$. Άρα, από το θεώρημα \eng{Bolzano}, υπάρχει τουλάχιστον ένα $x_0\in(0,1)$ τέτοιο ώστε
\[
\varphi(x_0)=0
\Leftrightarrow h(x_0)=e^{-x_0}.
\]

Για τη μοναδικότητα, για κάθε $x\in[0,1]$ είναι
\[
\varphi'(x)=\frac{1}{(x+1)^2}+e^{-x}>0.
\]
Άρα η $\varphi$ είναι γνησίως αύξουσα στο $[0,1]$, οπότε η ρίζα της είναι μοναδική.

Επομένως η εξίσωση
\[
h(x)=e^{-x}
\]
έχει ακριβώς μία ρίζα στο $(0,1)$.
\end{erwthma}
\end{thema}

\begin{thema}{Γ}
Δίνεται η συνεχής συνάρτηση $f:(0,+\infty)\to\mathbb{R}$ με
\[
f(e)=\frac{2}{e}
\]
και
\[
x^2f^2(x)-2x\ln x\cdot f(x)+\ln^2x-1=0.
\]
\begin{erwthma}
\item Η δοσμένη σχέση γράφεται
\[
\bigl(xf(x)-\ln x\bigr)^2-1=0
\Leftrightarrow \bigl(xf(x)-\ln x\bigr)^2=1.
\]
Θεωρούμε τη συνάρτηση
\[
F(x)=xf(x)-\ln x,\quad x>0.
\]
Η $F$ είναι συνεχής στο $(0,+\infty)$ και από την προηγούμενη σχέση ισχύει
\[
F(x)=1\quad\text{ή}\quad F(x)=-1.
\]
Επειδή το $(0,+\infty)$ είναι διάστημα και η $F$ είναι συνεχής, η $F$ δεν μπορεί να αλλάζει τιμή ανάμεσα στα $1$ και $-1$.

Επιπλέον
\[
F(e)=e f(e)-\ln e=e\cdot\frac{2}{e}-1=1.
\]
Άρα για κάθε $x>0$ είναι
\[
xf(x)-\ln x=1.
\]
Επομένως
\[
xf(x)=\ln x+1
\Leftrightarrow
{f(x)=\frac{\ln x+1}{x}},\quad x>0.
\]

\item Για κάθε $x>0$ είναι
\[
f'(x)=\left(\frac{\ln x+1}{x}\right)'
=\frac{\frac{1}{x}\cdot x-(\ln x+1)}{x^2}
=-\frac{\ln x}{x^2}.
\]
Βρίσκουμε ρίζες και πρόσημα της $f'$:
\begin{itemize}
\item $f'(x)=0\Leftrightarrow \ln x=0\Leftrightarrow x=1$.
\item $f'(x)>0\Leftrightarrow -\ln x>0\Leftrightarrow 0<x<1$.
\item $f'(x)<0\Leftrightarrow -\ln x<0\Leftrightarrow x>1$.
\end{itemize}
Στον παρακάτω πίνακα βλέπουμε τα πρόσημα της $f'$ και τη μονοτονία της $f$:
\begin{center}
\begin{tikzpicture}
\tkzTabInit[color,lgt=2,espcl=1.5,colorC = maincolor!50!white,colorL = maincolor!20!white,colorV = maincolor!50!white]%
{$x$ /0.8,$f'(x)$ /0.8,$f(x)$/0.8}%
{$0$,$1$,$+\infty$}%
\tkzTabLine{d, +, z, -, }
\tkzTabLine{d, \nearrow,t,\searrow,}
\end{tikzpicture}
\end{center}
Επομένως η $f$ είναι \textbf{γνησίως αύξουσα} στο $(0,1]$ και \textbf{γνησίως φθίνουσα} στο $[1,+\infty)$.

Άρα παρουσιάζει ολικό μέγιστο στο $x_0=1$, με τιμή
\[
f(1)=\frac{\ln1+1}{1}=1.
\]
Η $f$ δεν παρουσιάζει ολικό ελάχιστο, αφού
\[
\lim_{x\to0^+}\frac{\ln x+1}{x}=-\infty.
\]

\item Έχουμε
\[
\lim_{x\to+\infty}f(x)
=\lim_{x\to+\infty}\frac{\ln x+1}{x}
\xlongequal[\text{\eng{DLH}}]{\frac{+\infty}{+\infty}}
\lim_{x\to+\infty}\frac{\frac{1}{x}}{1}=0.
\]
Άρα η ευθεία
\[
{y=0}
\]
είναι οριζόντια ασύμπτωτη της $C_f$ στο $+\infty$.

\item Έστω $x(t)$ η τετμημένη και $y(t)$ η τεταγμένη του $M$ τη χρονική στιγμή $t$. Τότε
\[
y(t)=f(x(t))=\frac{\ln x(t)+1}{x(t)}.
\]
Το εμβαδόν δίνεται από τη σχέση
\[
E(t)=x(t)y(t)=x(t)f(x(t)).
\]
Όμως
\[
xf(x)=x\cdot\frac{\ln x+1}{x}=\ln x+1.
\]
Άρα
\[
E(t)=\ln x(t)+1.
\]
Παραγωγίζοντας ως προς $t$, παίρνουμε
\[
E'(t)=\frac{x'(t)}{x(t)}.
\]
Τη χρονική στιγμή $t_0$ είναι
\[
x(t_0)=e
\quad\text{και}\quad
x'(t_0)=2e^2.
\]
Επομένως
\[
E'(t_0)=\frac{2e^2}{e}=2e.
\]
Άρα ο ζητούμενος ρυθμός μεταβολής είναι
\[
{2e\ \text{μον}^2/\si{sec}}.
\]
\end{erwthma}
\end{thema}

\begin{thema}{Δ}
Θεωρούμε την παραγωγίσιμη συνάρτηση $f:(0,+\infty)\to\mathbb{R}$ για την οποία
\[
f(x)+e^{f(x)}=e^{x-1}-\ln x+\frac{e^{e^{x-1}}}{x}.
\]
\begin{erwthma}
\item Θέτουμε
\[
\varphi(x)=e^{x-1}-\ln x,\quad x>0.
\]
Τότε
\[
e^{\varphi(x)}=e^{e^{x-1}-\ln x}
=\frac{e^{e^{x-1}}}{e^{\ln x}}
=\frac{e^{e^{x-1}}}{x}.
\]
Άρα η δοσμένη σχέση γράφεται
\[
f(x)+e^{f(x)}=\varphi(x)+e^{\varphi(x)}.
\]
Θεωρούμε τη συνάρτηση
\[
g(y)=y+e^y,\quad y\in\mathbb{R}.
\]
Επειδή
\[
g'(y)=1+e^y>0
\]
για κάθε $y\in\mathbb{R}$, η $g$ είναι γνησίως αύξουσα, άρα είναι $1-1$.

Από τη σχέση
\[
g(f(x))=g(\varphi(x))
\]
και επειδή η $g$ είναι $1-1$, προκύπτει
\[
f(x)=\varphi(x)=e^{x-1}-\ln x.
\]
Επομένως
\[
{f(x)=e^{x-1}-\ln x,\quad x>0}.
\]

\item Για κάθε $x>0$ είναι
\[
f'(x)=e^{x-1}-\frac{1}{x}
\]
και
\[
f''(x)=e^{x-1}+\frac{1}{x^2}>0.
\]
Άρα η $f'$ είναι γνησίως αύξουσα στο $(0,+\infty)$.

Έστω $1<a<\beta$. Η $f$ είναι συνεχής στο $[a,\beta]$ και παραγωγίσιμη στο $(a,\beta)$, οπότε από το Θεώρημα Μέσης Τιμής υπάρχει $\xi\in(a,\beta)$ τέτοιο ώστε
\[
\frac{f(\beta)-f(a)}{\beta-a}=f'(\xi).
\]
Επειδή $\xi>a$ και η $f'$ είναι γνησίως αύξουσα, έχουμε
\[
f'(\xi)>f'(a).
\]
Άρα
\[
\frac{f(\beta)-f(a)}{\beta-a}>e^{a-1}-\frac{1}{a}.
\]

\item Έστω $x_0\in(1,2)$ το σημείο επαφής. Η εφαπτομένη της $C_f$ στο σημείο $A(x_0,f(x_0))$ έχει εξίσωση
\[
y-f(x_0)=f'(x_0)(x-x_0).
\]
Για να διέρχεται από το $O(0,0)$, πρέπει
\[
-f(x_0)=f'(x_0)(-x_0)
\Leftrightarrow f(x_0)=x_0f'(x_0).
\]
Άρα
\[
e^{x_0-1}-\ln x_0
=x_0\left(e^{x_0-1}-\frac{1}{x_0}\right)
\]
οπότε
\[
e^{x_0-1}(1-x_0)-\ln x_0+1=0.
\]
Θεωρούμε τη συνάρτηση
\[
\psi(x)=e^{x-1}(1-x)-\ln x+1,\quad x\in[1,2].
\]
Η $\psi$ είναι συνεχής στο $[1,2]$ και
\[
\psi(1)=1>0.
\]
Επίσης
\[
\psi(2)=-e-\ln2+1<0.
\]
Άρα, από το θεώρημα \eng{Bolzano}, υπάρχει τουλάχιστον ένα $x_0\in(1,2)$ τέτοιο ώστε
\[
\psi(x_0)=0.
\]

Για τη μοναδικότητα, για κάθε $x\in(1,2)$ είναι
\[
\psi'(x)=-xe^{x-1}-\frac{1}{x}<0.
\]
Άρα η $\psi$ είναι γνησίως φθίνουσα στο $(1,2)$, οπότε έχει το πολύ μία ρίζα.

Επομένως υπάρχει ακριβώς ένα $x_0\in(1,2)$ τέτοιο ώστε η εφαπτομένη της $C_f$ στο $A(x_0,f(x_0))$ να διέρχεται από την αρχή των αξόνων.

\item Για κάθε $x\in(1,e]$ ισχύει
\[
\ln x>0.
\]
Άρα
\[
f(x)=e^{x-1}-\ln x<e^{x-1}
\]
για κάθε $x\in(1,e]$. Επομένως
\[
\dint_1^e f(x)\,\d x<\dint_1^e e^{x-1}\,\d x.
\]
Υπολογίζουμε
\[
\dint_1^e e^{x-1}\,\d x=\left[e^{x-1}\right]_1^e=e^{e-1}-1.
\]
Άρα
\[
{\dint_1^e f(x)\,\d x<e^{e-1}-1}.
\]
\end{erwthma}
\end{thema}
%=========== ΤΕΛΟΣ ΛΥΣΗΣ - 2ο ΔΙΑΓΩΝΙΣΜΑ ΠΡΟΣΟΜΟΙΩΣΗΣ ===========
